\documentclass[11pt, a4paper]{article}

% --- Language & Font Setup ---
\usepackage{polyglossia}
\setmainlanguage{english}
\setotherlanguage{arabic}

% Use Noto Sans for English to ensure perfect support for ALA-LC 
% transliteration characters (macrons, underdots, and ayn 'ʿ').
\setmainfont{Noto Sans}

% Use Amiri for flawless Quranic Arabic rendering
\newfontfamily\arabicfont[Script=Arabic]{Amiri}

% --- Layout & Tables ---
\usepackage[top=2.5cm, bottom=2.5cm, left=1.2cm, right=1.2cm]{geometry}
\usepackage{booktabs}
\usepackage{array}
\usepackage{threeparttable} % Added for table notes

% Define custom column types for stacked Arabic + Transliteration
% 'm' vertically centers the cell contents relative to the row.
\newcolumntype{A}{>{\centering\arraybackslash}m{3.1cm}} % Main word columns
\newcolumntype{R}{>{\centering\arraybackslash}m{1.6cm}} % Root column
\newcolumntype{C}{>{\centering\arraybackslash}m{1cm}}   % Number column

\begin{document}
\thispagestyle{empty}
% \title{Arabic Numbers (1 -- 10)\\ \large Cardinal and Ordinal Forms with ALA-LC Transliteration}
% \author{}
% \date{}
% \maketitle


% \vspace{2em}

\begin{center}
	\begin{threeparttable}
		\renewcommand{\arraystretch}{1.6} % Add vertical padding to rows
		\begin{tabular}{ C R A A A A }
			\toprule
			                                                                         &   & \multicolumn{2}{c}{\textbf{Cardinal Numbers (1-10)}} & \multicolumn{2}{c}{\textbf{Ordinal Numbers (1st-10th)}} \\
			\cmidrule(lr){3-4} \cmidrule(lr){5-6}
			\textbf{Num}                                                             &
			\textbf{Root}                                                            &
			\textbf{w/ Masc. Noun} \tnote{(α)}                                       &
			\textbf{w/ Fem. Noun} \tnote{(α)}                                        &
			\textbf{Masculine}                                                       &
			\textbf{Feminine}                                                                                                                                                                             \\
			\midrule

			% 1
			{\Large\textarabic{١}}\newline {\large\textbf{1}}                        &
			{\Large\textarabic{و ح د}}\newline \textit{\small w-ḥ-d} \tnote{(β)}     &
			{\Large\textarabic{أَحَدٌ}}\newline \textit{\small aḥadun}                  &
			{\Large\textarabic{إِحْدَى}}\newline \textit{\small iḥdā}                   &
			{\Large\textarabic{الأَوَّلُ}}\newline \textit{\small al-awwalu}             &
			{\Large\textarabic{الأُولَى}}\newline \textit{\small al-ūlā}                                                                                                                                    \\
			\midrule

			% 2
			{\Large\textarabic{٢}}\newline {\large\textbf{2}}                        &
			{\Large\textarabic{ث ن ي}}\newline \textit{\small th-n-y}                &
			{\Large\textarabic{اِثْنَانِ}}\newline \textit{\small ithnāni}               &
			{\Large\textarabic{اِثْنَتَانِ}}\newline \textit{\small ithnatāni}            &
			{\Large\textarabic{ثَانٍ}}\newline \textit{\small thānin}                  &
			{\Large\textarabic{ثَانِيَةٌ}}\newline \textit{\small thāniyatun}                                                                                                                                 \\
			\midrule

			% 3
			{\Large\textarabic{٣}}\newline {\large\textbf{3}}                        &
			{\Large\textarabic{ث ل ث}}\newline \textit{\small th-l-th}               &
			{\Large\textarabic{ثَلَاثَةٌ}}\newline \textit{\small thalāthatun}           &
			{\Large\textarabic{ثَلَاثٌ}}\newline \textit{\small thalāthun}              &
			{\Large\textarabic{ثَالِثٌ}}\newline \textit{\small thālithun}              &
			{\Large\textarabic{ثَالِثَةٌ}}\newline \textit{\small thālithatun}                                                                                                                                \\
			\midrule

			% 4
			{\Large\textarabic{٤}}\newline {\large\textbf{4}}                        &
			{\Large\textarabic{ر ب ع}}\newline \textit{\small r-b-ʿ}                 &
			{\Large\textarabic{أَرْبَعَةٌ}}\newline \textit{\small arbaʿatun}             &
			{\Large\textarabic{أَرْبَعٌ}}\newline \textit{\small arbaʿun}                &
			{\Large\textarabic{رَابِعٌ}}\newline \textit{\small rābiʿun}                &
			{\Large\textarabic{رَابِعَةٌ}}\newline \textit{\small rābiʿatun}                                                                                                                                  \\
			\midrule

			% 5
			{\Large\textarabic{٥}}\newline {\large\textbf{5}}                        &
			{\Large\textarabic{خ م س}}\newline \textit{\small kh-m-s}                &
			{\Large\textarabic{خَمْسَةٌ}}\newline \textit{\small khamsatun}              &
			{\Large\textarabic{خَمْسٌ}}\newline \textit{\small khamsun}                 &
			{\Large\textarabic{خَامِسٌ}}\newline \textit{\small khāmisun}               &
			{\Large\textarabic{خَامِسَةٌ}}\newline \textit{\small khāmisatun}                                                                                                                                 \\
			\midrule

			% 6
			{\Large\textarabic{٦}}\newline {\large\textbf{6}}                        &
			{\Large\textarabic{س د س}}\newline \textit{\small s-d-s} \tnote{(γ)}     &
			{\Large\textarabic{سِتَّةٌ}}\newline \textit{\small sittatun}                &
			{\Large\textarabic{سِتٌّ}}\newline \textit{\small sittun}                   &
			{\Large\textarabic{سَادِسٌ}}\newline \textit{\small sādisun}                &
			{\Large\textarabic{سَادِسَةٌ}}\newline \textit{\small sādisatun}                                                                                                                                  \\
			\midrule

			% 7
			{\Large\textarabic{٧}}\newline {\large\textbf{7}}                        &
			{\Large\textarabic{س ب ع}}\newline \textit{\small s-b-ʿ}                 &
			{\Large\textarabic{سَبْعَةٌ}}\newline \textit{\small sabʿatun}               &
			{\Large\textarabic{سَبْعٌ}}\newline \textit{\small sabʿun}                  &
			{\Large\textarabic{سَابِعٌ}}\newline \textit{\small sābiʿun}                &
			{\Large\textarabic{سَابِعَةٌ}}\newline \textit{\small sābiʿatun}                                                                                                                                  \\
			\midrule

			% 8
			{\Large\textarabic{٨}}\newline {\large\textbf{8}}                        &
			{\Large\textarabic{ث م ن}}\newline \textit{\small th-m-n}                &
			{\Large\textarabic{ثَمَانِيَةٌ}}\newline \textit{\small thamāniyatun}         &
			{\Large\textarabic{ثَمَانٍ}}\newline \textit{\small thamānin} \tnote{(δ)}   &
			{\Large\textarabic{ثَامِنٌ}}\newline \textit{\small thāminun}               &
			{\Large\textarabic{ثَامِنَةٌ}}\newline \textit{\small thāminatun}                                                                                                                                 \\
			\midrule

			% 9
			{\Large\textarabic{٩}}\newline {\large\textbf{9}}                        &
			{\Large\textarabic{ت س ع}}\newline \textit{\small t-s-ʿ}                 &
			{\Large\textarabic{تِسْعَةٌ}}\newline \textit{\small tisʿatun}               &
			{\Large\textarabic{تِسْعٌ}}\newline \textit{\small tisʿun}                  &
			{\Large\textarabic{تَاسِعٌ}}\newline \textit{\small tāsiʿun}                &
			{\Large\textarabic{تَاسِعَةٌ}}\newline \textit{\small tāsiʿatun}                                                                                                                                  \\
			\midrule

			% 10
			{\Large\textarabic{١٠}}\newline {\large\textbf{10}}                      &
			{\Large\textarabic{ع ش ر}}\newline \textit{\small ʿ-sh-r}                &
			{\Large\textarabic{عَشَرَةٌ}}\newline \textit{\small ʿasharatun} \tnote{(ε)} &
			{\Large\textarabic{عَشْرٌ}}\newline \textit{\small ʿashrun} \tnote{(ε)}     &
			{\Large\textarabic{عَاشِرٌ}}\newline \textit{\small ʿāshirun}               &
			{\Large\textarabic{عَاشِرَةٌ}}\newline \textit{\small ʿāshiratun}                                                                                                                                 \\

			\bottomrule
		\end{tabular}

		\begin{tablenotes}
			\small
			\item[(α)] \textbf{Reverse Gender Agreement for some Ordinal numbers:} In Quranic/Classical Arabic, cardinal numbers 1 and 2 act as adjectives and agree with the noun's gender.
			However, cardinal numbers 3 through 10 exhibit \textit{reverse agreement} of gender.
			You must use the morphologically feminine number (ending in \textarabic{ة}) with masculine nouns, and the masculine number with feminine nouns.
			\item[(β)] \textbf{Roots for ``First'':} The ordinal numbers for first (\textit{awwal} / \textit{ūlā}) derive from the roots: \textarabic{أ و ل} (\textit{ʾ-w-l}).
			% \item \textbf{Roots for ``First'':} The cardinal numbers for one (\textit{aḥad} / \textit{wāḥid}) derive from the root \textarabic{و ح د} (\textit{w-ḥ-d}). However, the ordinal numbers for first (\textit{awwal} / \textit{ūlā}) derive from a completely different root: \textarabic{أ و ل} (\textit{ʾ-w-l}).
			\item[(γ)] \textbf{Number 6 Root:} The root for six is \textarabic{س د س} (\textit{s-d-s}).
			The \textit{dāl} assimilated into the \textit{tāʾ} over time, changing \textit{sidta} into \textit{sitta}.
			\item[(δ)] \textbf{Number 8:} When used with a feminine noun in the indefinite nominative/genitive state,
			the Classical standard dictates the deficient (manqūs) form \textarabic{ثَمَانٍ} (\textit{thamānin}), ending with kasratan.
			\item[(ε)] \textbf{Number 10:} Notice the difference in the medial consonant (\textarabic{ش}).
			It takes a \textit{fatḥah} (\textarabic{عَشَرَةٌ} \textit{ʿasharatun}) with masculine nouns,
			but a \textit{sukūn} (\textarabic{عَشْرٌ} \textit{ʿashrun}) with feminine nouns.
		\end{tablenotes}
	\end{threeparttable}
\end{center}

% % \vspace{1.5em}
% \noindent
% % \textbf{Grammar Note:} In Quranic/Classical Arabic, cardinal numbers 1 and 2 act as adjectives and agree with the noun's gender. However, cardinal numbers 3 through 10 exhibit \textit{reverse agreement} (polarity). You must use the morphologically feminine number (ending in \textarabic{ة}) with masculine nouns, and the masculine number with feminine nouns. 
% \textbf{Grammar Note:} In Quranic/Classical Arabic, cardinal numbers 1 and 2 act as adjectives and agree with the noun's gender. 
% However, cardinal numbers 3 through 10 exhibit \textit{reverse agreement} of gender. 
% You must use the morphologically feminine number (ending in \textarabic{ة}) with masculine nouns, and the masculine number with feminine nouns. 

\end{document}
